\needspace{9\baselineskip}
\begin{center}
\pdfbookmark[1]{3. On Art. 70.}{zusatz-003}
[3. On Art. 70.]
\end{center}
The angle \(N\) in the equations of Art. 62 is arbitrary, so put \(N=\lambda\); with \(r'=r\cos\beta\), \(R'=R\cos B\), \(\Delta'=\Delta\cos b\), these become:
\[
\begin{aligned}
1)\quad&r\cos\beta-R\cos B\cos(L-\lambda)=\Delta\cos b\cos(l-\lambda)\\
2)\quad&R\cos B\sin(\lambda-L)=\Delta\cos b\sin(l-\lambda)\\
3)\quad&r\sin\beta-R\sin B=\Delta\sin b.
\end{aligned}
\]

\SourceLeaf{302}From 2 and 3 we obtain
\[
\operatorname{tang}(l-\lambda)=\frac{R\cos B\sin(\lambda-L)}{r\cos\beta-R\cos B\cos(\lambda-L)}\,\mathpunct{\vcenter{\hbox{\rule{.45em}{.35pt}}}}
\]
or approximately
\[
4)\quad l-\lambda=\frac{R\cos B\sin(\lambda-L)}{r\cos\beta}\mathpunct{\cdot}
\]

Further, from 1 and 3 we obtain
\[
\begin{aligned}
&R\cos B\sin\beta\cos(\lambda-L)-R\sin B\cos\beta\\
&\qquad=\Delta(\sin b\cos\beta-\cos b\sin\beta\cos(l-\lambda))
\end{aligned}
\]
or
\[
\begin{aligned}
5)\quad&\frac{R\cos B\cos\beta(\operatorname{tang}\beta\cos(\lambda-L)-\operatorname{tang}B)}{r}\\
&\qquad=\frac{\Delta}{r}(\sin(b-\beta)+2\cos b\sin\beta\sin\tfrac12(l-\lambda)^2)
\end{aligned}
\]
or approximately
\[
=b-\beta.
\]
\begin{center}\rule{.18\textwidth}{.35pt}\end{center}
